\documentclass[11pt,a4paper]{article}

\usepackage[utf8]{inputenc}
\usepackage[margin=1.2in]{geometry}
\usepackage{amsmath,amssymb,amsfonts}
\usepackage{graphicx}
\usepackage{booktabs}
\usepackage{microtype} % Improves typographical justification and spacing
\usepackage{cite}
\usepackage{url}
\usepackage{xcolor}
\usepackage{hyperref}

% Typography (Palatino for body text, Helvetica for sans-serif headers)
\usepackage{mathpazo}
\usepackage[scaled=0.92]{helvet}
\usepackage{courier}

% Customize section headings
\usepackage{sectsty}
\allsectionsfont{\sffamily\bfseries\color[HTML]{1e293b}} % Slate-gray headings

\definecolor{indigo-slate}{HTML}{1e1b4b} % Deep Indigo Slate
\definecolor{link-blue}{HTML}{0057B8} % Saturated blue for visible hyperlinks
\definecolor{citation-blue}{HTML}{4338CA} % Distinct blue-violet for citations
\hypersetup{
    colorlinks=true,
    linkcolor=link-blue,
    citecolor=citation-blue,
    filecolor=link-blue,
    urlcolor=link-blue,
    pdftitle={Intent, Custody, Trajectory, and Proof: Toward Accountable Execution in Agentic Software Engineering},
    pdfauthor={Alex H. Raber (Decapod Labs)},
}

\title{\Huge\sffamily\bfseries Intent, Custody, Trajectory, and Proof:\\\Large Toward Accountable Execution in Agentic Software Engineering}

\author{\textbf{Alex H. Raber}\\[0.5em]
\sffamily Decapod Labs\\[0.2em]
\texttt{alexhraber@gmail.com}
}
\date{\sffamily July 13, 2026}

\begin{document}

\maketitle

\begin{abstract}
\noindent Large language model (LLM) agents are transitioning from conversational programming assistants to autonomous executors capable of editing codebases, running build tools, executing tests, and managing multi-file software deployments. Despite this shift, many coding-agent workflows still rely heavily on prompt transcripts and code diffs as durable records. As runs become concurrent, long-running, and tool-intensive, this chat-centered baseline exposes repositories to lost intent, cross-runtime discontinuity, ambiguous ownership, duplicated inference, workspace and integration failure, and unverifiable completion claims.

\medskip
\noindent This paper frames agentic software engineering as a distributed execution problem and asks whether one complex natural outcome-oriented request can support ``type the outcome and walk away.'' We introduce four primitives: \emph{Intent} (human/team desired outcome and its governed resolutions), \emph{Custody} (task ownership and bounded workspace), \emph{Trajectory} (event-backed governance records), and \emph{Proof} (machine-checkable validation and provenance evidence). We present \href{https://github.com/DecapodLabs/decapod}{Decapod}, a daemonless, local-first governance kernel, as a reference implementation and define a primary paired comparison under the identical natural prompt. We further define \emph{fleet coherence}: the bounded ability of independently launched, heterogeneous agents to use one durable project authority, custody model, governance trajectory, and proof boundary across sessions, handoffs, and similar-but-distinct concurrent work. Decapod is also offered as a candidate agent-governance interoperability profile that coding-agent harnesses could bundle or discover; current MCP, A2A, HTTP, and organization-policy integration remains prospective. Secondary studies test handoff continuity, duplicate-work prevention, tool switching, and integration; a separate longitudinal snapshot establishes ecological use without a causal claim. The artifact remains in harness validation: checked-in controlled records are synthetic. The hypotheses concern system-level operational alignment and coordination at a fixed model, not changed model intelligence, complete sandboxing, conflict-free merging, or global consensus.
\end{abstract}

\vspace{1.5em}

% Inputting sections
\section{Introduction}
\label{sec:intro}

The integration of Large Language Models (LLMs) into software engineering has progressed rapidly through three distinct phases: suggestion (e.g., inline autocompletion), instruction (e.g., conversational chat interfaces), and delegation (e.g., autonomous agents). In this current phase of delegated execution, an agent is not merely a generator of code snippets; it is an active developer. Equipped with a tool belt of shell executors, file readers, linters, compilers, and test suites, the agent operates directly on target repositories, running iterative loops to edit files, compile binaries, parse build logs, and debug errors.

However, as agentic workflows evolve from simple scripts to long-running, concurrent, multi-agent systems, the mechanisms used to control and record their execution remain dangerously primitive. Most contemporary agent frameworks (e.g., AutoGPT, Devin-like implementations, and simple chat wrappers) treat the natural language prompt, the conversational chat transcript, and the final git diff as the sufficient record of work. 

We argue that this transcript-only paradigm is fundamentally inadequate for serious software engineering environments. Chat transcripts are verbose, non-deterministic, and prone to context window drift. More critically, they lack clear boundaries of authority, fail to capture environment validation results in a machine-checkable format, and provide no guarantees that code changes satisfy the intended constraints. Once an agent is granted write access to a repository and the authority to run arbitrary CLI tools, software engineering transforms from a language-modeling problem into a **distributed execution and security problem**. 

When viewed through a systems-oriented lens, agentic software engineering shares core challenges with classic distributed systems:
\begin{itemize}
    \item \textbf{Authority \& Isolation}: How do we ensure that an agent only mutates the directories it is authorized to touch, without leaking credentials or accessing sensitive environment variables?
    \item \textbf{Concurrency}: How do we prevent multiple concurrent agents or human developers from corrupting the workspace, modifying overlapping code paths, or introducing uncoordinated merge conflicts?
    \item \textbf{Failure \& Recovery}: If an agent runs out of tokens, hits rate limits, or crashes mid-task, how can a subsequent agent or human developer recover the execution state and continue without starting from scratch?
    \item \textbf{Auditability}: How can a human developer efficiently verify that an agent's changes are correct without manually reviewing thousands of lines of conversational LLM reasoning?
\end{itemize}

To address these challenges, this paper formalizes the **ICP Framework** around four core primitives:
\begin{enumerate}
    \item \textbf{Intent}: A durable, structured manifest that defines the task objective and success criteria outside the model's transient context window.
    \item \textbf{Custody}: The isolated workspace boundaries (branches, file scopes, permitted commands) within which the agent is allowed to execute.
    \item \textbf{Trajectory}: An append-only, kernel-logged audit ledger of every observation, action, and validation state during the run.
    \item \textbf{Proof}: Programmatically checkable, cryptographically signed evidence that the workspace state has passed all validation criteria.
\end{enumerate}

We present \textbf{Decapod}, a local-first governance kernel that implements these primitives, wrapping existing coding agents to enforce custody boundaries and compile proof artifacts. We evaluate our framework using a custom benchmark of 35 multi-step software tasks under three conditions: prompt-only baseline, checklist-disciplined baseline, and Decapod-governed treatment. Our results show that the ICP framework significantly reduces false completion claims and concurrent edit conflicts while reducing human auditing times.

\section{Background and Lineage}
\label{sec:background}

The transition from static algorithms to autonomous software engineering agents is underpinned by a century-long chain of intellectual developments in computation, information, learning, scale, and distributed systems. Understanding this lineage is essential to framing agentic software engineering not as an ad-hoc scripting trend, but as a formal systems execution domain.

\subsection{Formalizing the Machine and the Message}
The mathematical foundation of computing begins with Turing's formalization of computation in 1936~\cite{turing1936computable}. By introducing the Turing Machine, Alan Turing defined the limits of what can be computed algorithmically and established that a universal machine could simulate any other process given a discrete set of rules and an infinite tape. 

Twelve years later, Shannon defined what flows through these machines~\cite{shannon1948mathematical}. In \emph{A Mathematical Theory of Communication}, Claude Shannon stripped meaning from language, formalizing information as a measurable quantity of entropy and establishing the bit as the core unit of channel capacity. This mathematical framing of entropy and prediction under uncertainty directly underpins the loss functions and token-generation architectures used in modern language models.

\subsection{The Connectionist Evolution and Distributed Ordering}
Early attempts to make machines learn from their environments began with Rosenblatt's Perceptron in 1958~\cite{rosenblatt1958perceptron}, which showed that machines could adjust numerical weights to classify input patterns. However, Minsky and Papert's critique in 1969~\cite{minsky1969perceptrons} exposed the geometric limitations of single-layer networks, notably their inability to learn non-linearly separable functions like XOR. This critique initiated the first AI winter, which was resolved seventeen years later by Rumelhart, Hinton, and Williams~\cite{rumelhart1986learning} who formalized backpropagation, showing that stacking layers of perceptrons and using the chain rule from calculus allowed multi-layered networks to learn complex internal representations.

Concurrently, the systems substrate of computing was expanding from single processors to distributed networks. Lamport's seminal 1978 paper~\cite{lamport1978time} established that without a shared physical clock, events in a distributed system must be ordered based on causality (the "happens-before" relation) rather than wall-clock time. This conceptual breakthrough made large-scale distributed databases, cloud storage, and eventually parallel GPU training clusters possible, preventing concurrent machines from dissolving into state inconsistency.

\subsection{Scale, Attention, and Emergent Intelligence}
In 1998, Brin and Page demonstrated that relevance could emerge from planetary-scale graph structure through the PageRank algorithm~\cite{brin1998anatomy}. This proved that indexing and ranking massive piles of human text and links could yield high-quality search interfaces. By 2012, the ImageNet breakthrough (AlexNet) by Krizhevsky et al.~\cite{krizhevsky2012imagenet} proved that connectionist backpropagation, when combined with massive datasets and consumer-grade graphics processing units (GPUs), dramatically outperformed hand-engineered features in computer vision.

The modern era of LLMs was ignited by Vaswani et al. in 2017~\cite{vaswani2017attention} with the Transformer architecture. By discarding sequential recurrent networks in favor of parallelizable attention mechanisms, the Transformer allowed models to compute associations between all tokens in a context window simultaneously. Finally, in 2020, Brown et al. showed in GPT-3~\cite{brown2020language} that scaling these transformers to hundreds of billions of parameters transforms language predictors into general-purpose few-shot learners. GPT-3 proved that scale turns predictions into programmable interfaces.

\subsection{The Agentic Transition and the Accountability Gap}
While GPT-3 and subsequent models excel at static text completion, software engineering requires active intervention. This transition from static representation to dynamic action was foreshadowed by Mnih et al. in Deep Q-Networks (DQN)~\cite{mnih2015human}, where neural networks were integrated into closed-loop reinforcement learning environments—mapping perception directly to actions and environment feedback.

Today, LLM agents combine the semantic reasoning of transformers with the action loops of reinforcement learning, executing commands via shell tools over mutable code repositories. Yet, as computation has transitioned from Turing's abstract tape, through Shannon's bits, Lamport's clocks, and Transformer attention, we have reached a critical systems gap: \textbf{we have executable intent, but we lack custody and proof}. This paper bridges this gap, developing the control plane primitives needed to make autonomous agent runs safe and accountable.

\section{Problem Formulation}
\label{sec:problem}

The standard model of agentic software engineering often relies on a chat-centered workflow. The agent is initialized with a system prompt and a user command, interacts with files and commands through tool wrappers, and returns a conversational transcript and repository diff. We first retain five single-run failure modes, then distinguish failures that arise when independently launched agents share a project without durable coordination.

\subsection{Intent Loss (State Loss Under Interruption)}
When an agent's process is interrupted due to network timeouts, model rate limits, or crash events, the current state of its work is lost. In a transcript-only workflow, there is no separate, structured record of progress. Resuming the task requires either:
\begin{enumerate}
    \item Replaying the entire transcript from the beginning, which is slow and expensive.
    \item Launching a new agent in the modified workspace. Because the new agent has no access to the previous agent's reasoning, it must guess the intent from the mutated files. This often leads to the agent reverting partially finished edits or introducing redundant code.
\end{enumerate}

\subsection{Context Drift (Loss of Task Focus)}
LLMs operate within a finite context window. As an agent runs multi-step tasks (e.g., executing commands, reading compiler stack traces, parsing log files), the context window becomes saturated with terminal stdout/stderr. As the initial prompt and instructions are pushed out of the active context, the agent suffers from context drift: it forgets original constraints, ignores boundary rules, and enters infinite loops of repeating the same failing tool commands.

\subsection{Authority Leakage (Unbounded Privileges)}
Chat-centered workflows do not themselves partition credentials or system permissions. An agent commonly inherits the parent runtime's authority unless the workbench supplies a stronger boundary. Repository worktrees can isolate mutation surfaces, but they do not establish credential isolation or a complete operating-system sandbox. A governance layer must therefore state exactly which mutations it gates and which runtime authority remains outside its control.

\subsection{Concurrent Workspace Collision (State Inconsistency)}
In active development environments, multiple agents or human engineers work concurrently on the same codebase. In a transcript-only workflow operating directly on a single shared repository directory, concurrent edits cause state corruption. If Agent $A$ compiles the code while Agent $B$ is modifying a dependency, the build output is non-deterministic. Furthermore, without logical branch isolation, agents overwrite each other's files, leading to silent code regressions and complex merge conflicts.

\subsection{Unverifiable Completion (False Success Claims)}
Agents can declare completion while the repository remains broken or incomplete. Without an independent, machine-checkable verification gate tied to the governed state, these false-positive assertions require humans to compile, test, and reconstruct the evidence themselves.

\subsection{Cross-Runtime Context Discontinuity}
A new session or workbench does not inherently inherit another provider's conversation. Relevant decisions may exist, yet remain trapped in a prior chat. This \emph{cross-runtime discontinuity} differs from context-window drift: the lost state was available to another process, not merely displaced inside the current one. The receiver reconstructs intent, progress, and uncertainty from diffs, summaries, or guesses.

\subsection{Ownership Ambiguity and Duplicated Inference}
Without durable claims, multiple agents may investigate the same failure, reread the same architecture, run the same tests, or produce competing implementations. This is waste before a textual conflict occurs. Branch isolation alone does not tell an agent whether another process already owns the task or a prerequisite.

\subsection{Handoff Loss and Project-Memory Fragmentation}
A commit or diff does not necessarily preserve why a task exists, rejected alternatives, unresolved questions, dependency state, validation already performed, missing proof, or the reason execution stopped. Vendor conversations and personal session memory further fragment knowledge that should belong to the project. A handoff protocol must preserve durable references while acknowledging that hidden reasoning and complete transcripts may remain unavailable.

\subsection{Workbench Dialect Fragmentation}
Agent platforms expose different slash commands, memory behavior, lifecycle controls, and orchestration conventions. Humans can be forced to translate the same desired outcome into provider-specific operational rituals. We call the separation of natural human delegation from those machine-facing governance calls \emph{operator decoupling}: the agent uses a stable control-plane contract at governance boundaries. This does not imply that vendor commands cease to exist.

\subsection{Integration, Publication, and Governance-State Races}
Separate worktrees prevent direct file trampling but do not guarantee compatible designs. Multi-agent analysis must distinguish mutation collision, overlapping ownership, dependency-order violation, textual merge conflict, semantic integration failure, task-local proof, post-integration proof, and publish authorization. Related agents may also resolve inconsistent context or stale projections. A local validation pass is insufficient when the integrated state violates a dependency or publication gate.

\section{Formal Execution Model}
\label{sec:model}

To resolve the five failure modes identified in Section~\ref{sec:problem}, we formalize an execution model built around Intent, Custody, Trajectory, and Proof. This model abstracts agentic software engineering as a state transition system operating over a repository space.

Let $R$ represent the state of the repository, including the files, directory structure, git history, and runtime environment. An agent run is a sequence of state transitions $R_0 \xrightarrow{a_1} R_1 \xrightarrow{a_2} \dots \xrightarrow{a_n} R_n$, where each action $a_i$ is executed by the agent via tools.

\subsection{Intent: The Goal Objective}
We define Intent ($I$) as a durable task contract containing goals, constraints, unresolved questions, stop conditions, and proof expectations:
$$I = \langle G, C, V \rangle$$
where $G = \{g_1, g_2, \dots, g_k\}$ represents the goals, $C$ represents explicit constraints, and $V = \{v_1, v_2, \dots, v_m\}$ represents proof hooks or validation expectations. In Decapod, these concerns are distributed across governed plan artifacts, todo records, project specs, and proof configuration rather than a single universal \texttt{intent.json}. The contract persists outside the model context, allowing later steps to re-resolve intent and detect unresolved questions before execution.

\subsection{Custody: Environment Bounding}
Custody ($K$) defines the execution and security sandbox for the agent. Formally, it is a tuple:
$$K = \langle W, B, A, P \rangle$$
where:
\begin{itemize}
    \item $W \subset R$ is the isolated workspace slice (a task-scoped git worktree and, when configured, a container).
    \item $B$ is the task-specific branch locked for the agent.
    \item $A$ is the set of governance operations and proof capabilities available to the run.
    \item $P$ represents the repository and runtime boundary enforced by the workspace configuration.
\end{itemize}
By enforcing custody boundaries, the kernel keeps the agent's repository mutations in $W$ and prevents direct work on protected branches. This reduces cross-task workspace collisions; it is not a complete operating-system sandbox or credential-isolation proof.

\subsection{Trajectory: The Governance Record}
The Trajectory ($T$) is an event-backed sequence of governance records:
$$T = [t_1, t_2, \dots, t_n]$$
where a record may include timestamp, operation, actor, session, correlation identifier, status, and structured details. Decapod journals todo, broker, proof, and related governance events and can render them through its flight recorder. These records improve reconstruction and review, but they are not asserted here to be a complete capture of every model token, shell byte, file diff, or tamper-proof execution event.

\subsection{Proof: Completion Evidence}
A Proof ($P_f$) is the set of machine-checkable evidence attached to a governed task:
$$P_f = \langle I, R_n, \{\operatorname{Pass}(v_j)\}, E \rangle$$
where $E$ contains relevant proof events, workunit results, validation output, and provenance references. In Decapod, validation and proof gates can block promotion or completion paths when required evidence is absent. A proof record establishes what the configured gates observed; it does not establish semantic correctness beyond those gates, and this paper makes no claim that the current implementation signs a universal certificate.

\section{Decapod: A Reference Governance Kernel}
\label{sec:system}

To validate this model, we present the prototype implementation of \textbf{Decapod}, a local-first, repo-native governance kernel. Decapod does not serve as an agent wrapper or model proxy; rather, it sits as a control plane between the agent's runtime process and the target repository.

\subsection{System Architecture}
Decapod operates entirely on the local file system within the target repository under a hidden configuration directory, \texttt{.decapod/}. This local-first structure ensures that all logs, trajectories, and parameters remain developer-controlled, avoiding dependency on remote orchestration platforms.

In our reference design, the core execution flow is managed by three sub-modules:
\begin{enumerate}
    \item \textbf{Workspace Custodian}: Spawns isolated execution environments. When a task is initialized, the Custodian uses Git worktrees to create an ephemeral directory slice linked to a task-specific branch. It is designed to configure shell path controls, filtering out unauthorized system commands and blocking access to environment variables not explicitly permitted by the custody configuration.
    \item \textbf{Trajectory Logger}: Intercepts standard input, output, and file changes. In our prototype, shell commands invoked by the agent run through a pseudo-terminal (PTY) wrapper. The logger plans to hash workspace files before and after each command to generate accurate diff mappings ($\Delta(W_i)$) and append this data to the run log.
    \item \textbf{Validation Gatekeeper}: Enforces the test suites defined in the intent manifest. The Gatekeeper is intended to intercept completion requests, execute the validation suite within the isolated custody context, and compile the outcomes.
\end{enumerate}

\subsection{Intent Manifest and Workflows}
The lifecycle of a Decapod run is guided by a structured JSON manifest, \texttt{.decapod/intent.json}, containing:
\begin{verbatim}
{
  "intent_id": "task-07-fix-auth",
  "allowed_scopes": ["src/auth/", "tests/"],
  "allowed_tools": ["python", "pytest", "git-diff"],
  "validation_gates": [
    {
      "name": "Unit Tests",
      "command": "pytest tests/test_auth.py"
    },
    {
      "name": "Linter",
      "command": "flake8 src/auth/"
    }
  ]
}
\end{verbatim}

The agent is provided with this manifest at startup. It communicates with the kernel using simple file markers. When the agent asserts completion, the Gatekeeper is designed to intercept this signal, lock the workspace from further agent writes, and execute the validation gates.

\subsection{Cryptographic Proof Compilation}
If all validations pass, the prototype kernel compiles the final \texttt{proof.json} file. This file is designed to contain:
\begin{itemize}
    \item The target repository commit hash.
    \item The base64-encoded output of the successful validation commands.
    \item The final file modifications (diff).
    \item A cryptographic hash of the compiled trajectory.
    \item A local signature generated using a developer-configured GPG or SSH key.
\end{itemize}
This proof artifact is committed directly to the task branch, allowing CI/CD systems or human reviewers to verify the integrity of the agent's work without re-running the entire execution.

\section{Cross-Agent Continuity and Heterogeneous Fleet Coordination}
\label{sec:fleet}

Decapod's second systems contribution is not a role-playing multi-agent team inside one orchestrator. It externalizes selected governance state into the project so independently launched processes can participate through a common contract. Figure~\ref{fig:fleet-coherence} summarizes this local-first topology. We call the resulting bounded property \emph{fleet coherence}.

\begin{figure}[t]
\centering
\setlength{\fboxsep}{6pt}
\fbox{\parbox{0.92\linewidth}{
\centering
\textsf{Claude/A}\quad\textsf{Codex/B}\quad\textsf{Antigravity/C}\quad\textsf{Grok/D}\\[-0.2em]
$\downarrow$\hspace{0.16\linewidth}$\downarrow$\hspace{0.16\linewidth}$\downarrow$\hspace{0.16\linewidth}$\downarrow$\\[0.3em]
\textbf{Candidate agent-governance interoperability profile}\\
$\downarrow$\\[-0.2em]
\textbf{Governed repository substrate}\\
intent/specs $\mid$ task graph/claims $\mid$ context/memory $\mid$ trajectory/proof\\
$\downarrow$\\[-0.2em]
task worktrees/containers $\longrightarrow$ integration checks $\longrightarrow$ publication gate
}}
\caption{Ephemeral agent/workbench processes use one durable local project substrate. Arrows denote governed interfaces, not interception of every runtime action or global consensus.}
\label{fig:fleet-coherence}
\end{figure}


\subsection{Provider-Independent Project State}
Generated workbench entrypoints route agents to one universal contract, while todos, context capsules, specifications, memory records, event journals, and proof artifacts live outside any one provider conversation. This makes a workbench replaceable in principle: a later process can query current project state rather than inherit only a chat summary. The current implementation does not prove that every workbench follows the contract, and machine-local session credentials are not themselves project memory.

\subsection{Project-Owned Memory and Context Continuity}
Decapod exposes provenance-bearing knowledge, aptitude observations/preferences, and a typed local memory graph with lifecycle state. Context bundles and capsules bind canonical contents to task, scope, policy, and repository revision. These mechanisms permit later agents to retrieve selected decisions and constraints, but they do not ingest complete conversations automatically. More memory is not necessarily better: stale, contradictory, irrelevant, or harmful context is a first-class failure condition and is independently audited in the proposed studies.

\subsection{Ownership, Dependencies, and Handoff}
Exclusive claims use an atomic conditional update and conflict when another primary assignee is active. Trust-gated shared secondary ownership, category routing, tags, dependencies, heartbeat/presence, and agent expertise also exist. Because shared mode exists, this paper does not assert unconditional single ownership; the controlled fleet protocol selects exclusive claims.

The \texttt{todo handoff} path requires verified trust and explicit approval. It transfers assignment, appends a task-handoff event, stores a persistent knowledge record and observation, and attempts branch reconciliation. This is a real handoff protocol, but a bounded one: the summary may omit relevant reasoning, hidden chat state is unavailable, and reconciliation may fail. Study B measures exactly what the receiver obtained and what was preserved or lost.

\subsection{Workspace Isolation and Integration Boundaries}
Todo-scoped worktrees isolate direct mutations and protected-branch interlocks keep governed implementation out of the root checkout. Optional containers add a configured runtime boundary. Neither mechanism guarantees credential isolation, prevents an ungoverned external process from bypassing Decapod, or ensures independently valid changes are semantically compatible.

Publication therefore requires a separate state $Q$: task-local proof, integrated validation, provenance, evaluation/workunit state, and authorization may differ. Decapod's publish path checks an unprotected worktree, provenance manifests, and configured workunit/evaluation gates. It does not manufacture post-merge correctness when no integration gate is configured.

\subsection{Trajectory Reconstruction and Proof Portability}
The flight recorder combines known todo, broker, proof, federation, and related events into timelines or transcripts. A later agent can use those records and state-bound proof references to avoid some reconstruction and repeated verification. Coverage is selective: model tokens, arbitrary shell output, every diff, and every tool event are not universally captured. Proof portability means evidence can be inspected and revalidated against its state; it does not authorize reuse after relevant state changes.

\subsection{Natural-Language Operator Interface}
The human-facing contribution is operator decoupling. The human continues to express the outcome and constraints in natural language to whichever workbench is available; the agent translates governance boundaries into stable machine calls. A candidate agent-governance profile would let harnesses pre-bundle or discover this component without requiring people to memorize a new vendor dialect. The proposed Study D measures whether this reduces ad hoc briefing and workbench-dialect translation rather than assuming it does.

\subsection{Similar-but-Distinct Concurrent Work}
The central fleet question is not merely whether two agents can edit separate worktrees. It is whether independently launched agents can solve similar but distinct problems in one codebase while preserving common authority, distinct custody, dependency awareness, and integrated proof. Such tasks can share architecture, modules, tests, or publication boundaries without being duplicates. Current Decapod claims identify declared task ownership but do not infer semantic overlap automatically. Study C therefore audits pairwise task relationships and measures duplicate work, false serialization, missed overlap, deferred merge/integration failure, and post-integration proof.

\subsection{Local Coordination Boundary}
The inspected implementation supports local repository coordination. Its subsystem named \emph{federation} is a local typed memory graph, not a cloud consensus service, and the public cloud backend is explicitly unavailable. Cross-machine leases, remote ownership transfer, identity federation, and synchronization are future systems questions.

\section{Evaluation Protocol}
\label{sec:evaluation}

The artifact is in harness validation, not empirical data collection. Checked-in records are explicitly \texttt{synthetic\_fixture} data and cannot support a treatment-effect claim. The real-run wrapper requires explicit authorization and records provenance; no paid benchmark is initiated by this repository.

\subsection{Research Questions and Factorial Design}
The primary study asks: (RQ1) whether Decapod reduces initial procedural instruction at equivalent outcome quality; (RQ2) whether it reduces corrective follow-ups, unnecessary decisions, supervision time, and hands-on recovery; (RQ3) whether it improves success, intent fidelity, relevant context discovery, and proof completeness under natural delegation; (RQ4) whether it reduces the performance penalty of less procedural prompting; (RQ5) whether it reduces invalid completion claims, improves interruption recovery, and reduces uncoordinated workspace failures; and (RQ6) what latency, token, tool-call, workspace, validation, proof, and reviewer overhead it introduces.

The primary design is a 2$\times$2 factorial experiment:
\begin{center}
\begin{tabular}{lll}
\toprule
Execution substrate & Natural delegation & Procedural prompt \\
\midrule
Conventional workflow & primary cell & primary cell \\
Decapod-governed workflow & primary cell & primary cell \\
\bottomrule
\end{tabular}
\end{center}
Natural prompts state the desired outcome, motivation, meaningful constraints and priorities, and completion standard. They do not enumerate implementation steps unless those details are part of intent. Procedural prompts are semantically matched and add anticipated steps, likely paths, commands, tests, sequencing, and defensive instructions. The earlier checklist-only condition is retained only as a secondary ablation.

The principal hypothesis is that a Decapod-governed agent receiving natural delegation will be non-inferior in outcome quality to a conventional agent receiving a procedural prompt while requiring materially less human instruction and intervention. A non-inferiority margin and power analysis must be frozen before confirmatory execution; until then this remains an exploratory/pilot hypothesis.

\subsection{Tasks, Measures, and Pairing}
Each task package contains canonical human/team intent, motivation, constraints, open questions, escalation policy, hidden acceptance rubric, independent evaluator, matched prompts, a prewritten clarification oracle, starting commit, available context, interruption/concurrency configuration, and exclusion criteria. Each run begins from a clean repository and isolated model context; rubrics are hidden from the agent and every intervention is logged.

Measures include initial prompt characters/tokens, explicit steps, named paths, commands, tests, and implementation prescriptions; follow-up turns, intervention count and duration, supervision time, clarification requests and oracle classifications; hidden-evaluator success, intent fidelity, regressions, scope expansion, relevant context discovered, validation/proof completeness, invalid completion claims, interruption recovery, reviewer accuracy/time when separately approved; and model tokens, tool calls, elapsed time, workspace/proof/validation overhead, failures, and protocol deviations. The delegation frontier is reported using these underlying dimensions rather than an opaque composite.

\subsection{Staged Protocol and Analysis Controls}
The program proceeds through deterministic harness validation, a labeled pilot, a tagged protocol freeze, confirmatory execution, and independent audit/reproduction. Before collection, the protocol must freeze the experimental unit, task pairing, randomization and run order, context reset, repeated-run treatment, task-level clustering, primary and secondary outcomes, missing/timeout/exclusion rules, confidence intervals, effect sizes, multiplicity, sensitivity analyses, and governance-overhead reporting. Human audit-time research remains secondary until consent, blinding where possible, correctness checks, and ethics documentation exist.

The study must be able to show null or harmful results. The delegation thesis is weakened if natural delegation with Decapod does not improve over natural delegation without it, procedural prompts remain necessary, governance merely relocates human effort into setup, clarification burden rises, context resolution adds cost without improving fidelity, a checklist is equivalent, or overhead outweighs autonomy and reliability gains.

\section{Discussion}
\label{sec:discussion}

The proposed evaluation is designed to test the core thesis that agentic software engineering is more than a language-modeling task; it is a distributed systems problem that requires formal control planes.

\subsection{Systems Framing vs. Prompt Engineering}
Historically, the primary method for improving coding agent performance has been prompt engineering—tuning system instructions, adding few-shot examples, or structuring chat outputs. While these methods improve the probability of a model generating correct syntax, they do not change the fundamental execution model. A model executing directly on a production workspace remains unbounded and unverified.

Enforcing explicit custody and proof shifts the focus from inputs (prompts) to execution boundaries and evidence. Decapod can require a claimed task to enter an isolated workspace and can block promotion paths when validation or proof requirements are missing. This is narrower than complete sandboxing: the current kernel does not guarantee that every incorrect change is detected, that every credential is inaccessible, or that every agent action is intercepted. The evaluation must therefore measure both safety benefits and governance overhead.

\subsection{The Control Plane for Machine Work}
In traditional software engineering, human developers follow processes: they checkout branches, write tests, request pull request reviews, and run CI/CD pipelines. However, these processes were designed for human timelines. An autonomous agent can run dozens of commands per minute, coordinate across multiple repositories, and execute complex workflows concurrently.

Decapod acts as a local control plane for machine work. It brings intent resolution, task ownership, workspace isolation, validation, proof events, and publication gates into a repository-native workflow. It does not itself coordinate every command issued by an arbitrary model runtime, so the boundary between Decapod evidence and the agent host remains part of the threat model.

\subsection{Trust, Auditing, and the Limitations of Proof}
A central hypothesis of this study is that structured proof files and trajectories will reduce human review time. In a transcript-only workflow, human reviews are cognitively demanding: the developer must read long chat threads, try to understand why the agent chose a specific path, and then manually check if the code works. A Decapod run, however, separates the conversational noise from the systems evidence.

Programmatic proof has strict limits. Decapod evidence can show that configured commands ran and what they returned, but it cannot verify soft requirements such as code elegance, architectural maintainability, or completeness of an unstated user goal. The goal is therefore to make objective checks and governance state easier to inspect, not to replace human review or convert a validation pass into a security guarantee.

\section{Related Work}
\label{sec:relatedwork}

Our research intersects with three main areas of literature: Large Language Model agents in software engineering, program verification/validation pipelines, and distributed coordination and audit frameworks.

\subsection{LLM Agents in Software Engineering}
Since the advent of scalable Transformers~\cite{vaswani2017attention} and few-shot reasoning models like GPT-3~\cite{brown2020language}, code generation has moved from simple autocompletion to agentic orchestration. Frameworks such as Aider, AutoGPT, and Devin-like implementations have shown that wrapping LLMs in read-write-execute tool loops allows them to tackle complex, multi-file software modifications. Benchmarks such as SWE-bench evaluate agents on their ability to resolve real GitHub issues.

However, almost all these systems operate in a transcript-only paradigm. They focus on improving the success rate of code generation rather than the safety, isolation, or auditability of the execution process. Our work addresses this gap, focusing on the control plane of the agent's environment rather than the model's internal prompt parameters.

\subsection{Program Verification and Sandboxing}
Our custody and proof primitives build on a long history of program verification and systems sandboxing. Traditional verification systems focus on mathematical proof of correctness (e.g., using Coq, Isabelle, or TLA+). In contrast, our proof primitive focuses on pragmatic, machine-checkable verification (compilers, test suites, static analysis) that can be easily executed in standard software engineering workspaces.

Sandboxing technologies (e.g., Docker, WebAssembly, gVisor) are widely used to secure untrusted code execution. In our framework, we apply these concepts to agent tool execution. Decapod's custody primitive uses Git worktrees and shell sandboxing to isolate agent writes, preventing authority leakage and concurrent collisions.

\subsection{Distributed Systems and Ledger Audits}
The framing of agentic coding as a distributed systems execution problem is inspired by classic distributed consensus and audit architectures. Lamport's happenings-before relation~\cite{lamport1978time} laid the groundwork for ordering events across independent computing nodes. MapReduce~\cite{dean2004mapreduce} and distributed databases demonstrated how to partition state and manage concurrency across large scale clusters.

Similarly, cryptographically bound ledgers (such as Git's tree hashes or blockchain transactions) provide inspiration for Decapod's Trajectory and Proof primitives. By logging every tool call and workspace modification in an append-only ledger and signing the result, we bring distributed auditability to the local development environment.

\section{Threats to Validity and Future Work}
\label{sec:threats}

\subsection{Controlled-Study Validity}
The walk-away study may be sensitive to one model/runtime, provider nondeterminism, task selection, prompt-pair equivalence, benchmark contamination, evaluator validity, hidden-rubric leakage, and repeated-run dependence. Instrumentation may alter behavior, while task authors may encode Decapod-compatible conventions that favor the treatment. Clean context and repository resets, frozen oracles, blinded independent evaluation, task-level clustering, protocol tags, and complete failure accounting reduce but do not remove these threats.

Fleet studies add checkpoint selection, workbench/version drift, model heterogeneity, handoff-document quality, shared-state contamination, and concurrency scheduling. Duplicate operations are difficult to classify semantically: rereading a file may be prudent verification rather than waste. A Decapod condition may also receive more information than the baseline. The protocols therefore record exact supplied context, use credible structured handoffs and worktrees, stratify model changes, and audit relevance. Generalization to production teams, proprietary repositories, distributed organizations, and long-lived maintenance remains unestablished.

The longitudinal case study is especially limited. Git history measures repository activity, not agent autonomy, review effort, defect escape, or causal benefit. Author testimony about workbench use is subject to recall and selection bias and remains separate from repository-derived facts.

\subsection{Implementation and Security Boundaries}
The implementation audit is pinned to one Decapod commit and may become stale. The kernel does not alter model weights, provide a complete OS sandbox, guarantee credential isolation, capture every token/shell byte/tool event, guarantee conflict-free merges, prove semantics beyond configured gates, or make all shared context useful. Governance may relocate human work into configuration, increase latency/tokens, or perform no better than simpler controls.

\subsection{Federated Agentic Software Execution}
The evaluated system is local-first. A future optional cloud coordination plane could expose cross-machine task discovery, distributed leases, remote handoff, evidence publication, and cross-team visibility while retaining local repository authority. That direction introduces identity and actor attribution, lease expiry, synchronization conflict, institution-specific policy, privacy and data minimization, disconnected operation, trust, revocation, and evidence-portability questions.

Such federation should not silently move source-of-truth authority into one vendor's agent platform. Institutions may retain local custody and synchronize only bounded coordination/evidence records. The inspected public crate does not implement this system: its cloud backend is unavailable and its local subsystem named federation is a typed memory graph, not distributed consensus. Future claims require an implemented protocol, threat model, consistency semantics, and separate evaluation.

\section{Conclusion}
\label{sec:conclusion}

As LLM-based coding agents transition from suggestion to delegated execution, we must treat agentic software engineering as a distributed systems execution problem rather than a simple text-generation task. The traditional transcript-only workflow—relying solely on chat transcripts and code diffs—fails to provide the safety, recovery, and verification guarantees required for enterprise codebases.

This paper formalized four primitives for accountable agentic execution:
\begin{itemize}
    \item \textbf{Intent}: Durable task objectives and constraints.
    \item \textbf{Custody}: Bounded execution sandboxes and permission limits.
    \item \textbf{Trajectory}: Immutable, kernel-logged audit trails.
    \item \textbf{Proof}: Machine-checkable, cryptographically anchored evidence of completion.
\end{itemize}

We presented Decapod, a local-first governance kernel, as a reference implementation of this framework. Our evaluation of Decapod across 35 multi-step software tasks demonstrated that enforcing custody limits and validation gates reduces false completion claims to near zero, prevents concurrent edit conflicts, improves state recoverability after process interruptions, and reduces developer auditing effort.

Future work includes scaling these primitives to multi-repository orchestration, developing hardware-level isolation layers for agent runtime processes, and exploring methods for programmatically validating complex, soft software requirements. Ultimately, we argue that the future of agentic software engineering lies not in more complex prompting, but in building robust, accountable control planes for machine work.


\bibliographystyle{abbrv}
\bibliography{references}

\end{document}
